\chapter{The Sentence the Interior Cannot Settle}

The boundary returned a bit and left a claim unnamed. A balance fails to close because the
apparatus can state something it cannot settle from within the account that states it; the bit
was the mark of that failure, read from outside. This chapter names the claim.

It is a sentence the apparatus writes in its own account. Through every prior chapter the
apparatus read things it could carry --- a mark, a receipt, an energy, a residue. Here it
carries a sentence about its own closure: a claim, written in the order the account already
speaks, that says the account closes in that order. The sentence is admissible. The apparatus
can hold it, encode it, and set about settling it. The one thing it cannot do is settle it from
inside.

Every route the apparatus has runs through the interior, and every interior route, carried to
its end, arrives at the boundary the last chapter read --- the balance that will not close, the
residue that will not vanish. The sentence is precisely the claim whose attempted settlement is
that residue. There is no shorter road and no other road: the interior, asked to decide the
sentence, returns the obstruction every time.

So the sentence is true, or it is false, and the apparatus cannot say which from inside. This is
not a failure of effort, and not a limit of size. It is the fact the whole movement has been
approaching --- a claim the account can frame and cannot close, whose standing is decided not by
any interior route but by the single mark the boundary returns. The interior writes the
sentence; the outside reads its bit.

The chapter reads this in order: the claim the interior carries; its reduction, by every route,
to one finite obstruction; the standing only the outside distinguishes; the recurrence of the
sentence one step further on; and, last and most carefully, what the chapter does and does not
claim --- for the sentence wears the shape of a famous one, and the shape is all it wears.

\section{The Claim the Interior Carries}

The previous chapter ended with a mark and an unnamed claim. The mark was not large enough to
tell a story by itself. It said only that the boundary balance did not close. But a balance does
not fail to close in the abstract. It fails in answer to something the apparatus has tried to
settle. Chapter 19 begins by naming that something. The claim is not brought in from outside the
account to explain the bit after the fact. It is a sentence the apparatus can write in the same
order in which it has carried every earlier closure.

That order is the closure order. It is the book's native way of saying that one settled piece
stands at or below another: one receipt under a later receipt, one account under the next account,
one closed stage under the stage that is supposed to contain it. The order is not an extra
language placed above the apparatus. It is how the apparatus already keeps track of its own
stability. When a closure stands under another closure, the apparatus can say so. When a source
closure is assigned a target closure one step on, the apparatus can hold the pair and ask whether
the standing really holds.

The sentence is exactly that assertion of standing. It names a source closure, names a target
closure, and says that the source stands at or below the target in the account's own order. The
content is deliberately spare:
\[
  C_{\mathrm{source}} \leq C_{\mathrm{target}} .
\]
The notation is only a reader's handle for the claim, not a new theory hovering behind it. It
says that the account closes as the account says it closes. The source is what the apparatus
already carries; the target is the next closure the same order authorizes; the sentence asserts
their relation.

This makes the claim project-native. It is not a borrowed sentence wearing the apparatus's
clothes. It is not a statement about all formal systems, not a report about a hidden semantics,
and not a claim that some outside theory has already decided. It is born from the apparatus's
own bookkeeping: closure below closure, source toward target, held relation in the order the
account itself speaks. If the earlier chapters taught the apparatus to carry marks, receipts,
energies, residues, and signs, this chapter asks it to carry one more kind of object: a sentence
about the closure of that carrying.

The sentence is therefore admissible, but admissibility must stay modest. To say that the
apparatus can write the sentence is not to say that it can settle the sentence. Writing means the
claim has a place in the account. Holding means the account can keep it fixed long enough to ask
what follows from it. Encoding means it can be carried as a finite object rather than as a vague
intention. Carrying means it can be passed through the account's routes and tested against the
closures those routes provide. None of those acts is a decision. They make the claim available
to the apparatus. They do not make the claim resolved by the apparatus.

This is the point at which the chapter has to resist an easy overreading. A sentence about the
account is not automatically a sentence the account can judge. The apparatus has spoken about
itself before, but self-reference here is not a magical device and not a license to import a
famous outside result. The claim is internal because it uses the closure order. It is self-directed
because it concerns the account's own standing. It is admissible because it can be written,
held, encoded, and carried. Those are finite permissions. They are not yet settlement.

The distinction matters because the previous chapter already showed what happens at the end of
an attempted settlement. The boundary balance does not close; the residue remains; the outside
reading returns a mark. If admissibility were the same as settlement, the mark would have no
work left to do. The apparatus would write the claim and thereby own its decision. But Chapter
18 forbade that shortcut. The mark appears precisely because the account can arrive at a
boundary with no interior term left to spend. So the sentence must be strong enough to be a real
claim in the account and weak enough, by itself, not to decide its own standing.

Its strength is that it is about closure, the very relation the apparatus has used to organize
its prior readings. It does not describe an external world, predict an observation, or reach for
a quantity beyond the boundary. It asks whether the account stands in the relation it says it
stands in. Its weakness is that the same closure order provides only the routes of the interior.
The apparatus can take the sentence along those routes, but it cannot step outside those routes
while remaining the interior apparatus. That is where the obstruction from Chapter 18 will
return.

So S01 names the claim but does not try to finish it. The sentence is a carried assertion of
closure: source under target, account under its own next account, standing stated in the account's
own order. It is not a world claim, not a measurement claim, and not a universal theorem claim.
It is the precise sort of sentence the apparatus is entitled to frame. The next section asks
what happens when the apparatus tries to settle it. The answer is not a decision from inside,
but a reduction to the finite obstruction already waiting at the boundary.

\section{The Reduction to a Finite Obstruction}

To settle the sentence the apparatus does the only thing it can do from inside: it carries the
claim through its own account. The sentence is admissible, so the carrying can begin. It has a
source closure, a target closure, and a stated relation between them. The apparatus asks whether
that relation can be closed by the routes already available to the account. No new exterior
voice is added at this point, and no wider theory is invited to decide the matter. The apparatus
uses its own closure order and follows the claim as far as that order permits.

The carrying is finite. That is the first guardrail on the section. The attempted settlement is
not an endless wandering through possible interpretations, and it is not a loop in which the
apparatus keeps rephrasing the same claim without arriving anywhere. The claim is written in a
finite account, carried by finite terms, and tested against finite closures. Each route the
interior owns has a terminus. The question is not whether the apparatus can search forever. It
cannot, and it does not need to. The question is what the finite carrying reaches when it has
spent the routes available to it.

It reaches the boundary already read in Chapter 18. The source closure is carried toward the
target closure; the account asks whether the stated standing closes; the first-order terms are
given their chance to answer. They close as far as they can close. The shaped boundary balance is
then exposed again: demand at the boundary, carried part of the account, resistance, passage,
and no interior change left to use. Nothing in the sentence adds a hidden interior correction.
Nothing in the attempt supplies a term that Chapter 18 had not already counted. So the carrying
comes to the same place: the balance that does not close.

Write the attempted settlement as the residue it produces. Carried to its end, the sentence does
not return an interior yes or an interior no. It returns the obstruction:
\[
  \text{settle}(\text{sentence}) \;=\; \operatorname{res} \;=\; -1 \;\neq\; 0 .
\]
The notation should be read with the same modesty as the previous section's closure notation. It
is not a machinery for deciding all statements, and it is not a hidden predicate of truth. It is
a compact way to say what happens to this claim in this apparatus. The attempted settlement
terminates, and the thing it terminates at is nonzero.

That nonzero terminus is the whole reduction. The sentence is not undecidable because the
apparatus has failed to be clever enough, and not because the apparatus has been stopped before
doing its work. It is the opposite: the apparatus has done its work and arrived at a definite
finite object. The object is the residue, the unit already carried as the single invariant and
read at the boundary as non-closure. A decision would have closed the route. The residue leaves
the route terminated without decision.

This distinction matters. An interior failure can sound like silence, as if the apparatus asks
its question and hears nothing back. That is not what happens here. The interior is not silent.
It answers with the obstruction. It can display the route, carry the sentence, reach the
boundary, and name the residue. What it cannot do is turn that residue into a decision about the
sentence from within the same account. The answer it owns is exactly the wrong kind of answer
for settlement: not true, not false, but nonzero obstruction.

Nor is the obstruction a second claim smuggled into the argument. It is the same finite remainder
Chapter 18 isolated in the outside reading. There the residue made the boundary signal on. Here
the same residue is what the attempted settlement of the carried sentence reaches. The chapter
therefore does not need to appeal to a classical result about formal systems, and it does not
need a special predicate that says when a sentence is provable. Its reduction is project-native:
carry the closure claim; exhaust the interior closures; arrive at the non-closing boundary
balance; read the residue that remains.

The reduction also explains why the sentence from S01 was worth naming. A claim about closure
can be carried to the exact place where closure fails. If the sentence had been about something
outside the account, the route would have depended on an imported subject matter. If it had been
only a loose metaphor for self-reference, the route would have had no finite terminus. Instead
the sentence concerns the account's own closure order, and that order carries it to the boundary
where the account has already learned to read non-closure. The obstruction belongs to the
sentence because the attempted settlement of the sentence produces it.

So ``the interior cannot settle the sentence'' is not a vague confession. It is a located fact:
the apparatus reaches \(\operatorname{res}=-1\neq0\) when it tries to settle the claim. The
residue is definite, but it is not a decision. The route terminates, but it terminates at the
mark of non-closure. That is why the next section can ask a sharper question. If the interior
always reaches the same obstruction, then the two standings of the sentence are not
distinguished from inside. Whatever distinguishes them must be read where Chapter 18 learned to
read it: at the boundary, by the single returned bit.

\section{Distinguished Only from Outside}

From inside, the sentence's two standings are not kept apart. The account can write the claim,
carry the claim, and reach the obstruction produced by the claim. It can also write the opposed
standing and carry that opposed standing through the same closures. But when either standing is
handled as an interior matter, the carrying arrives at the same finite remainder. The account
does not find one internal route that closes and another that refuses to close. It finds one
shape of failure: the route terminates at the boundary with the residue still present.

That sameness is not a defect in notation. It is the point of the section. An interior
distinction must be made by something the interior owns: a term, a closure, a comparison, a
finite event that separates one standing from the other before the boundary is reached. Here
there is no such separator. The sentence is about the account's own closure, and the account's
own closure is exactly where the residue appears. The claim and its opposed standing therefore
do not split inside the apparatus. They travel to the same place and leave the same mark of
non-closure behind.

The apparatus can see that something has gone wrong with settlement. It can name the residue. It
can say that the attempted decision has not become an interior decision. Those are genuine
interior achievements. But none of them tells the apparatus which standing of the sentence has
been distinguished. The interior knows the obstruction as obstruction. It does not own the act
that reads the obstruction as the standing of the sentence. That act occurs at the boundary,
where Chapter 18 located the returned bit.

The boundary bit is correspondingly narrow. It does not open a window onto everything the
apparatus fails to know, and it does not replace the apparatus with a wiser judge. It says only
that the obstruction is present. In the language of the previous chapter, the bit is on when the
boundary balance does not close. Here that same on-reading attaches to the carried sentence:
the finite attempt to settle the sentence has reached \(\operatorname{res}=-1\neq0\), so the
outside reads the obstruction as present. The mark distinguishes the standing from outside
because the interior has already exhausted its own routes. No extra content is added by the
mark; the mark only receives the exhausted route as a completed boundary fact.

This produces a two-level result. At the first level, the account has no interior distinction
between the two standings. Both are carried to the same nonzero residue. At the second level,
the boundary reading distinguishes the fact of non-closure. It does not hand the interior a new
route; it does not make the sentence internally settled after all. It records, from outside, the
finite fact that the interior route failed to close. The sentence is therefore not decided
inside and then reported outside. It is undecided inside, and the outside reads that
undecidedness by the one mark the boundary can return.

The word ``outside'' can mislead if it is heard as distance or privilege. The outside is not a
larger theory hovering above the apparatus. It is the position of reading the boundary result
after the apparatus has completed its own finite carrying. From that position, the residue is no
longer a task to be settled by the same account. It is a returned fact. The outside does not
have to rerun the interior's closures, because the closures have already reached their terminus.
It only reads the terminus as a boundary signal.

So the difference between the two standings is not hidden somewhere inside the sentence. It is
not a missing clause, a cleverer encoding, or a longer carrying. If such an interior separator
existed, the account would have to display it as part of its own closure order. The reduction in
the previous section says that no such separator is produced: every owned route reaches the
same nonzero obstruction. The difference appears only when the boundary result is read as a
result, rather than treated as one more interior task.

This is why the chapter can speak of a sentence the interior cannot settle without turning that
sentence into a mystery. The failure is fully located. The account can show the route of the
failure, and the outside can read the mark returned by that route. What the account cannot do is
convert the mark into an interior separation of the sentence's standings. Inside, the standings
collapse into the same residue. Outside, the residue is readable as the on-bit of non-closure.

The distinction also protects the sentence from becoming too large. The chapter is not saying
that every difficult claim gains an outside answer, or that every boundary mark decides a world.
It is saying something smaller and sharper: this carried closure-sentence produces a finite
obstruction when the interior tries to settle it, and that obstruction is distinguished only by
the boundary reading. The mark reads presence of residue, not a hidden catalogue of all truths.
The account remains the account; the outside remains the boundary reading of what the account
has already returned.

\section{The Sentence Recurs}

The sentence is not a single accident of the account. Once the apparatus can frame a closure
claim at one stage, it can frame the next closure claim by the same finite act that gives it a
successor anywhere else in the account. Nothing exotic is added. The apparatus has already
learned how to pass from a stage to its next stage; it uses that passage here. One step further
along the closure order stands another claim of the same kind: that the account closes, now one
closure higher than before.

The next claim is not a new subject matter. It repeats the same project-native form at the next
level. The apparatus asks again whether the account's own closure order can settle the claim it
has just carried. The carrying is again finite. The available routes are again the routes owned
by the account. And the route again reaches the same boundary obstruction rather than an
interior decision. The first sentence did not exhaust a special trick; it exposed a pattern in
the way the closure order treats claims about its own closure.

This is the reader-facing force of the successor-of-successor idea. The account does not need
to prove counting from the beginning each time a new sentence is framed. It already owns the
finite successor step, and that step can be applied to the stage at which the previous closure
claim stood. From the first carried sentence, the apparatus can step to the next carried
sentence. From that one, it can step again. Each step gives a new sentence of the same type:
framed inside, carried inside, not settled inside.

The recurrence is therefore exact but modest. It does not say that every possible statement has
been sorted into a grand hierarchy. It says that the account's own finite successor operation
keeps producing closure-claims that return to the same kind of boundary test. At each stage the
interior can frame the claim and spend its closures. At each stage the attempted settlement
terminates at the nonzero residue. At each stage the boundary bit reads the presence of that
residue from outside.

So the limit the chapter found is not a single edge. It is a standing feature of the closure
order. The first sentence shows the edge once; the next sentence shows that the edge recurs
when the account advances one step. The apparatus does not run out of such sentences by
completing the first one, because the operation that produced the first remains available for
the next. Nor does the recurrence make the outside reading broader than it was. The bit remains
the same narrow mark: obstruction present, closure not completed, no interior settlement
returned.

\section{What Is and Is Not Claimed}

The sentence wears a familiar shape, and the chapter must say plainly that the shape is all it
wears. It is a closure-claim the account can frame about its own closure, and it is not settled
by the routes the account owns from inside. That resemblance is useful because it tells the
reader where the pressure lies. It is also dangerous if it is allowed to carry more than the
construction has earned.

The positive claim is finite and local. This particular carried sentence, in this particular
account, reaches no interior settlement. When the account spends the closures available to it,
the attempted settlement ends with the same nonzero obstruction isolated earlier in the chapter.
From outside, the boundary bit reads that obstruction as present. That is the whole affirmative
content: one framed closure-sentence, one finite obstruction, one outside mark that records
non-closure. The claim is not portable by slogan. It depends on the actual finite route the
chapter has followed: frame the sentence, carry it through the account's closure order, meet
the residue, and read the residue at the boundary.

Nothing stronger follows from this chapter. It does not prove a classical incompleteness
theorem. It does not introduce an arithmetized provability predicate, a truth predicate, or a
semantic oracle. It does not assert a theorem about all formal accounts, all possible
sentences, or a general world state. It does not say that physical measurement, by itself,
decides mathematical truth. It also does not prove a continuum theorem, an existence theorem,
or a smoothness theorem. Those are different claims, with different burdens.

The point is instead narrower and more mechanical. The account forms enough interior language
to carry this sentence. The same account lacks an interior route that distinguishes its two
standings once the finite closure attempt has been made. The outside bit does distinguish them,
but only as a boundary reading of residue: obstruction present rather than obstruction absent.
It is not a hidden judge of every sentence. It is not a replacement semantics. It is a mark
attached to the failure of this finite closure. The chapter therefore does not ask the reader
to trust a broad analogy. It asks the reader to follow a bounded calculation to the point where
the inside has no remaining separating move.

This is why the chapter can use the old shape without borrowing the old theorem. The shape says:
the account can form a claim about its own settling. The construction says: in this instance,
the settling attempt returns a nonzero remainder. The boundary says: that remainder is visible
from outside. The chapter claims exactly that much and no more. Its conclusion is a located
boundary fact, not a universal law.

\section*{Bridge: True from False}

The chapter began with a sentence the account could carry but could not settle. That was the
first important narrowing. The sentence was not imported as a grand theorem about every
possible formal language. It was made inside the account's own finite furniture: closure,
stage, successor, residue, and boundary reading. The apparatus could state the closure-claim in
its own terms. It could move the claim through the routes it owns. It could spend the available
closures. What it could not do was finish the settling from inside.

The obstruction was also kept finite. The chapter did not appeal to a hidden authority beyond
the calculation. It followed the carried sentence to the place where the attempted interior
settlement leaves a nonzero remainder. That remainder is the obstruction Chapter 19 has been
using all along. It is not a mood, not a mystery, and not a general declaration that interior
accounts are doomed. It is the residue returned by this finite attempt. When the account asks
for closure on the sentence that says it has closure, the route does not return closure. It
returns a boundary obstruction.

The outside-only distinction is the second narrowing. Inside, the two standings of the
sentence are not separated by a route the account can use. The same carried form reaches the
same interior failure of separation. From outside, however, the boundary bit can read the
residue as present. That outside mark does not say everything true about the sentence. It does
not give a full semantics for the account. It says only that the finite obstruction is there.
The bit is narrow enough to remain honest: obstruction present, closure not completed,
interior settlement not returned.

The recurrence makes the limitation stable without making it universal. Once the account has
the successor step that moves from one closure stage to the next, the same kind of
closure-claim can be framed one step higher. The next sentence is not a new metaphysical
creature. It is the same project-native pressure applied at the next stage. It asks again
whether the account's own closure order can settle the claim it has carried. It again reaches
the nonzero obstruction. The edge therefore recurs as the account advances, but it recurs by a
finite step. The chapter has not climbed to a theorem about all sentences; it has shown that
this boundary test can be repeated by the account's own successor operation.

This recurrence matters for the reader's path through the book. Earlier chapters taught the
apparatus to distinguish, count, carry, and read finite remainders. Chapter 19 uses those same
habits on a sentence about closure itself. The result is not a decorative self-reference. It is
the moment when the account's ordinary finite operations are aimed back at the account and find
their own boundary. That is why the next chapter can begin with a bit rather than with a new
axiom. The bit is already the record left by the attempt.

The claim gate then fixed the scope. Chapter 19 claims this: a carried closure-sentence,
framed inside the finite account, reaches no interior settlement and leaves the nonzero
obstruction that the outside bit reads. It does not claim a classical incompleteness theorem.
It does not introduce a provability predicate, a truth predicate, or an all-purpose oracle. It
does not decide a general world state, settle physical measurement as a theorem, or prove a
continuum, existence, or smoothness result. The resemblance to older self-reference is a shape,
not a borrowed result. The chapter earns only the located boundary fact it has traced.

That leaves a new problem, and the bridge must not solve it too early. Chapter 19 has mostly
used the bit in its on position. The obstruction is present; the closure attempt has not
completed; the outside reading marks that fact. But a bit has two values. If the on value marks
non-closure, then the off value must have its own role. The account now has to say what it is
for the boundary to read no residue, and it has to do so without pretending that the interior
has suddenly learned the distinction on its own.

The bridge also changes the reader's attention. Up to this point the question was whether the
sentence could be settled inside. Now the question becomes what distinction the boundary
reading can support once the inside has failed. The chapter does not answer that by naming a
finished semantics. It passes forward a simpler task: compare the obstruction-present case
with the obstruction-absent case and see whether the boundary can keep them apart.

This is where true and false enter as the next boundary separation. The interior route did not
hold them apart for the carried sentence. It could frame the sentence and fail to settle it,
but it could not produce the separating mark. The outside bit is the first place where the
difference has the right form: on is not off. Chapter 20 takes up that difference. It will
read the obstruction-present case against the no-obstruction case and ask what sort of truth
and falsity can be named at the boundary when the interior route has collapsed them.

So the handoff is exact. Chapter 19 supplies the sentence, the finite obstruction, the
outside-only mark, the recurrence, and the scope limit. Chapter 20 receives a two-valued
boundary reading whose values still need to be distinguished. The next step is not to make the
outside bit wiser than it is. The next step is to watch the bit do the one thing the interior
could not do here: separate true from false at the boundary.
